\chapter*{Liste des acronymes et abréviations}
\addcontentsline{toc}{chapter}{Liste des acronymes et abréviations}

\begin{longtable}{@{}lp{11cm}@{}}
\textbf{AUC} & \textit{Area Under the Curve}: aire sous la courbe ROC utilisée ici comme mesure descriptive exploratoire \\[3pt]
\textbf{CSV} & \textit{Comma-Separated Values}: format de fichier tabulaire \\[3pt]
\textbf{CUSUM} & \textit{Cumulative Sum}: somme cumulative pour la détection de changements persistants \\[3pt]
\textbf{FAIR} & \textit{Findable, Accessible, Interoperable, Reusable}: principes de gestion des données scientifiques \\[3pt]
\textbf{IF} & \textit{Isolation Forest}: algorithme de détection d'anomalies par isolation \\[3pt]
\textbf{HYPO} & Branche temporelle recherchant une baisse persistante et multivariée de l'activité \\[3pt]
\textbf{INSTABILITÉ} & Branche temporelle recherchant un mouvement disproportionné, des transitions denses et une posture fragmentée \\[3pt]
\textbf{IoU} & \textit{Intersection over Union}: mesure de recouvrement temporel \\[3pt]
\textbf{JSON} & \textit{JavaScript Object Notation}: format d'échange de données \\[3pt]
\textbf{LOF} & \textit{Local Outlier Factor}: détecteur d'anomalies fondé sur la densité locale \\[3pt]
\textbf{MAD} & \textit{Median Absolute Deviation}: écart absolu médian \\[3pt]
\textbf{MI} & \textit{Motion Index}: indice de mouvement fourni par le capteur \\[3pt]
\textbf{PLF} & \textit{Precision Livestock Farming}: élevage de précision \\[3pt]
\textbf{SHA-256} & Fonction de hachage cryptographique utilisée pour vérifier l'intégrité des artefacts \\[3pt]
\textbf{SLS} & Score observationnel de la cohorte IceTag--SLS, somme de quatre composantes binaires dans les fichiers étudiés \\[3pt]
\end{longtable}

\newpage

\noindent\textbf{Termes techniques}

\vspace{0.5em}

\begin{tabular}{@{}lp{10cm}@{}}
\textbf{Alerte comportementale} & Signal indiquant une dégradation persistante de plusieurs familles de mesures; elle doit être vérifiée et n'est pas un diagnostic \\[3pt]
\textbf{Période de référence} & Partie initiale de la série d'une vache utilisée pour estimer son comportement attendu \\[3pt]
\textbf{Boiterie} & Altération de la locomotion dont l'évaluation clinique repose sur la démarche, la posture, l'appui et, au besoin, l'examen du pied \\[3pt]
\textbf{Couverture} & Proportion des mesures attendues effectivement observées dans un intervalle \\[3pt]
\textbf{Créneau horaire} & Position dans le cycle journalier utilisée pour comparer un signal à des périodes homologues de la période de référence \\[3pt]
\textbf{Événement en période future} & Perturbation injectée uniquement après la partie utilisée pour construire la référence \\[3pt]
\textbf{Inspection visuelle} & Contrôle humain de la cohérence technique des événements et des sorties, sans adjudication clinique \\[3pt]
\textbf{Notification de fond} & Notification observée hors fenêtre injectée sur des données non étiquetées; elle ne peut être qualifiée de vraie ou fausse alerte clinique \\[3pt]
\textbf{Notification de vérification} & Sortie de priorité 2 ou 3 demandant une observation ciblée de l'animal, sans conclure à une boiterie \\[3pt]
\textbf{Période réfractaire} & Durée minimale entre deux notifications; elle dépend de la branche et vaut 12 heures pour la fusion principale \\[3pt]
\textbf{Chaîne de traitement (pipeline)} & Suite des opérations allant de l'ingestion des données aux épisodes et aux notifications \\[3pt]
\end{tabular}
